\section*{Chapitre \ref{ch:g6:decimals} --- Les nombres décimaux}

\begin{solution}{exo:g6:decimals:1}
$5.37$~; \quad $0.9$~; \quad $12.004$~; \quad $8.05$.
\end{solution}

\begin{solution}{exo:g6:decimals:2}
Chiffre des dixièmes~: $3$~; chiffre des centièmes~: $5$~; le $8$
compte les dizaines.
\[
86.354 = 8 \times 10 + 6 + 3 \times \frac{1}{10} + 5 \times
\frac{1}{100} + 4 \times \frac{1}{1000}.
\]
\end{solution}

\begin{solution}{exo:g6:decimals:3}
$5.3 > 5.29$ (comparer $5.30$ avec $5.29$)~; \quad
$0.7 = 0.70$ (les \omterm{def:g1:counting:zero}{zéros} finaux ne changent rien)~; \quad
$2.09 < 2.9$ (dixièmes~: $0 < 9$)~; \quad
$14.5 > 14.49$.
\end{solution}

\begin{solution}{exo:g6:decimals:4}
Compléter à deux décimales~: $4.20$~; $4.05$~; $4.51$~; $4.15$~;
$4.50$. Ordre~:
\[
4.05 < 4.15 < 4.2 < 4.5 < 4.51 .
\]
\end{solution}

\begin{solution}{exo:g6:decimals:5}
Chaque graduation est un dixième. $A$~: $6.3$~; \quad $B$~: $6.7$~;
\quad $C$~: $7.2$.
\end{solution}

\begin{solution}{exo:g6:decimals:6}
$45.80 + 7.65 = 53.45$.

$23.40 - 8.72 = 14.68$ (vérif.~: $14.68 + 8.72 = 23.4$).

$56 \times 24 = 1344$, et deux chiffres décimaux dans les facteurs~:
$5.6 \times 2.4 = 13.44$.
\end{solution}

\begin{solution}{exo:g6:decimals:7}
$6.42 \times 10 = 64.2$~; \quad
$0.35 \times 1000 = 350$~; \quad
$78.1 \div 100 = 0.781$~; \quad
$5 \div 1000 = 0.005$.
\end{solution}

\begin{solution}{exo:g6:decimals:8}
$2.4$~m $= 240$~cm~; \quad $370$~g $= 0.37$~kg (diviser par $1000$)~;
\quad $0.85$~km $= 850$~m~; \quad $56$~mm $= 0.056$~m.
\end{solution}

\begin{solution}{exo:g6:decimals:9}
À l'unité~: $24$ (chiffre suivant $8$~: \omterm{def:g4:numbers:round}{arrondir} vers le haut). Au
dixième~: $23.9$. Au centième~: $23.87$. Et $9.97$ au dixième~: le
chiffre suivant est $7$, donc on arrondit vers le haut les $9$
dixièmes --- $10.0$.
\end{solution}

\begin{solution}{exo:g6:decimals:10}
Prix des baguettes~: $3 \times 1.15 = 3.45$. Monnaie~:
$5 - 3.45 = 1.55$.
\end{solution}

\begin{solution}{exo:g6:decimals:11}
Entre $7.4$ et $7.5$~: par exemple $7.45$ (ou $7.41$, $7.499$,
\dots). Entre $3.99$ et $4$~: par exemple $3.995$. Dans les deux cas
il y a \emph{infiniment} de tels nombres~: on peut toujours ajouter
des décimales.
\end{solution}

\begin{solution}{exo:g6:decimals:12}
Le plus grand~: mettre les plus grands chiffres en premier et la
\omterm{def:g4:decimals:point}{virgule} le plus tard possible~: $85.2$. Le plus petit~: les plus
petits chiffres en premier et la \omterm{def:g4:decimals:point}{virgule} le plus tôt possible~:
$2.58$.
\end{solution}

\begin{solution}{pb:g6:decimals:1}
\textbf{1.} $3$ est le plus grand~: en complétant, $3.000 > 2.999$.
La \omterm{ex:g1:subtraction:difference}{différence} vaut $3.000 - 2.999 = 0.001$, un millième.

\textbf{2.} Strictement entre $7.4 = 7.40$ et $7.5 = 7.50$ se
trouvent
\[
7.41,\ 7.42,\ 7.43,\ 7.44,\ 7.45,\ 7.46,\ 7.47,\ 7.48,\ 7.49 :
\]
neuf nombres, un par nouvelle graduation.

\textbf{3.} La même image, dix fois plus petite~: entre $7.440$ et
$7.450$ se trouvent $7.441, 7.442, \dots, 7.449$ --- neuf nombres à
nouveau.

\textbf{4.} Entre $7.400$ et $7.500$, les nombres à trois chiffres
après la \omterm{def:g4:decimals:point}{virgule} sont $7.401, 7.402, \dots, 7.499$~: tous les
millièmes de $401$ à $499$, soit $99$ nombres.

\textbf{5.} La recette du zoom~: écrire les deux nombres avec le même
nombre de décimales, puis \emph{ajouter une décimale de plus} --- le
plus petit des deux suivi d'un chiffre de $1$ à $9$ tombe strictement
entre les deux. En effet $5.67 = 5.670$ et $5.68 = 5.680$, et
\[
5.670 < 5.675 < 5.680 ;
\]
de même $0.1999 = 0.19990 < 0.19995 < 0.20000 = 0.2$. Entre deux
graduations voisines, il y a toujours tout un nouvel étage de neuf
graduations plus fines.

\textbf{6.} $3 < 3.05 < 3.1$~; puis $3 < 3.005 < 3.01$~; puis
$3 < 3.0005 < 3.001$. Chaque candidat est battu par un zoom de plus.

\textbf{7.} Quel que soit le nombre proposé par Tom --- son candidat,
strictement plus grand que $3$ --- la question 5 fabrique un nombre
strictement compris entre $3$ et ce candidat. Le candidat n'était
donc \emph{pas} le nombre le plus proche de $3$~: il en existe un
encore plus proche. Comme cela arrive à tous les candidats sans
exception, aucun nombre ne peut être «~le nombre juste après $3$~»~:
il n'existe tout simplement pas.

\textbf{8.} $2.9 < 2.95 < 3$, \ $2.99 < 2.995 < 3$, \
$2.999 < 2.9995 < 3$. Et $3 - 2.999 = 0.001$. Avec vingt neufs, la
\omterm{ex:g1:subtraction:difference}{différence} est une \omterm{def:g4:decimals:point}{virgule} suivie de dix-neuf \omterm{def:g1:counting:zero}{zéros} et d'un $1$ ---
une unité de la vingtième décimale~: minuscule, mais pas nulle. Quel
que soit le nombre de neufs que Tom écrit, il n'atteint jamais $3$.

\textbf{9.} Les nombres entiers montent par pas de $1$~: être plus
grand que $7$ et plus petit que $8$ demanderait un nombre entier
d'unités strictement compris entre $7$ et $8$ unités, et il n'y en a
pas. Le zoom n'y échappe qu'en écrivant des chiffres \emph{après la
\omterm{def:g4:decimals:point}{virgule}} --- précisément ce que les nombres entiers n'ont pas.

\textbf{10.} La plus petite longueur qui s'arrondit en $3.7$ est
$3.65$ cm~: son chiffre suivant est $5$, donc on arrondit \emph{vers
le haut} (\cref{def:g6:decimals:rounding}). Et $3.75$ s'arrondit en
$3.8$ pour la même raison. Ainsi «~$3.7$ cm au dixième près~»
signifie~: au moins $3.65$ cm, et strictement moins que $3.75$ cm ---
tout un segment de longueurs possibles se cache derrière une seule
valeur arrondie.

\textbf{11.} $0.0999 < 0.123 < 0.58 < 0.6$. En une phrase~: en
comparant chiffre par chiffre depuis la gauche, $0.0999$ a $0$
dixième alors que $0.6$ en a $6$, et la comparaison est réglée dès là
--- le nombre de chiffres écrits ne dit rien de la taille
(\cref{met:g6:decimals:compare}).

\textbf{12.} Non~: en centimes, $4.99$ euros font $499$ centimes et
$5$ euros font $500$ centimes --- deux nombres \emph{entiers}
consécutifs, sans rien entre eux. Les prix, qui ont exactement deux
décimales, sont en réalité des nombres entiers de centimes
déguisés~: comme les nombres entiers de la question 9, ils ont bel et
bien des voisins de palier. Le zoom sans fin exige le droit d'écrire
toujours plus de décimales.

\textbf{13.} Les candidats proches de $6$ sont $5.82$ et $8.25$ (un
nombre commençant par $2$, ou par $58$, $25$, $82$, $85$, est loin de
$6$). Distances~: $6 - 5.82 = 0.18$ et $8.25 - 6 = 2.25$. Le plus
proche est $5.82$.

\textbf{14.} Par exemple $3.1415$, puisque
$3.1410 < 3.1415 < 3.1420$~; puis $3.14159$, puisque
$3.14150 < 3.14159 < 3.14160$. (Ce sont deux étapes réelles de la
chasse à $\pi = 3.14159\dots$)

\textbf{15.} Supposons que quelqu'un nomme un nombre, aussi grand
qu'il le veut. Un zoom transforme chaque couple de graduations
voisines entre $0$ et $1$ en neuf nouveaux nombres, et le zoom peut
se répéter indéfiniment --- les questions 2, 3 et 4 n'étaient que les
deux premiers étages. En le répétant assez de fois, on \omterm{def:g2:mult:def}{produit} plus
de nombres décimaux entre $0$ et $1$ que le nombre annoncé. Aucun
nombre n'est donc assez grand pour les compter~: il y en a une
infinité.
\end{solution}
